% Expose IX, sections 11.4 and 11.5.
% French transcription lines 3148--3191; authority scan physical pp. 466--468
% (scan indices 465--467; source folios 454--456). Physical p. 468 was
% checked through the exact 11.6 boundary. All pages were compared in four
% overlapping 1100-dpi bands. The text, formulas, primes, checks, underlines,
% and displayed arrows are clear at that resolution; no escalation was needed.

\medskip
\noindent
11.4. The preceding description of a pairing (11.3.1), by means of the
expression (11.2.1) for the $\ell$-primary component of the group of
connected components, is valid only when $\ell\ne p$. Suppose now that
$A_K$ has \underline{semistable reduction} over $S$. We shall again
define pairings that are perfect dualities

\begin{equation}\tag{11.4.1}
\Phi_o(\ell)\otimes\Phi'_o(\ell)\longrightarrow\mathbb Q/\mathbb Z
\end{equation}
without any restriction on $\ell$. As in 11.3,\footnote{The printed
source reads ``as in 11.4'' here, within section 11.4 itself. The
reference is corrected to the immediately preceding construction in
11.3.} one should establish that this recovers (up to a sign, perhaps,
which ought to be specified) the pairing defined in (1.2.1). This would
therefore establish Conjecture 1.3 in the case of semistable reduction,
just as it does when one restricts to the $\ell$-primary components for
$\ell\ne p$.

To define (11.4.1), we shall give, essentially, an expression for
$\Phi_o(\ell)$ in terms of the Barsotti--Tate group $T_\ell(A_K)$ over
$K$; in the present case this will play the same role as 11.2 did in
the case treated above.

\medskip
\noindent
\underline{Theorem} 11.5. \underline{Suppose that $S$ is henselian and
that $A_K$ has semistable reduction over $S$. Let $\ell$ be a prime and
consider the monodromy pairing} (9.1.2), \underline{and hence the
canonical homomorphism}

\begin{equation*}
\widetilde u_\ell:\underline M_\ell
 \longrightarrow\underline{\check M}'_\ell.
\end{equation*}
\underline{Put $\Phi(\ell)=\operatorname{coker}\widetilde u_\ell$.
Then there is a canonical isomorphism}

\begin{equation}\tag{11.5.1}
\Phi_o(\ell)\xrightarrow{\ \sim\ }\Phi(\ell)\times_S s.
\end{equation}

This statement immediately gives the definition of a perfect pairing
(11.4.1). Indeed, applying 11.5 to $A'_K$ in place of $A_K$ and using
(10.2.2), we also obtain a canonical isomorphism

\begin{equation*}
\Phi'_o(\ell)\xrightarrow{\ \sim\ }\Phi'(\ell)\times_S s,
\qquad
\Phi'(\ell)=\operatorname{coker}\bigl(
 {}^t\!\widetilde u_\ell:\underline M'_\ell
 \longrightarrow\underline{\check M}_\ell\bigr).
\end{equation*}

\medskip
\noindent
\underline{Remarks} 11.5.2. a) Theorem 11.5 implies in particular that
$\widetilde u_\ell$ is an isogeny, that is, that $u_\ell$ is
separating, including in the case $\ell=p$, where this assertion is not
immediate from the definition in 9.6. Notice that we do not use here
the (considerably deeper) Theorem 10.4, which will be proved in the next
section.

\smallskip
\noindent
b) Using 10.4, which supplies us with a monodromy pairing

\begin{equation*}
u:\underline M\otimes\underline M'\longrightarrow\mathbb Z_S,
\end{equation*}
and denoting by $\Phi$ the cokernel of the corresponding homomorphism

\begin{equation*}
\widetilde u:\underline M\longrightarrow\underline{\check M}',
\end{equation*}
we obtain an equivalent formulation of 11.5, independent of the choice
of $\ell$, in the canonical isomorphism

\begin{equation}\tag{11.5.3}
\Phi_o\xrightarrow{\ \sim\ }\Phi\times_S s.
\end{equation}
